% ==============================================================================
% SECCIÓN 11: GESTIÓN DE CALIDAD Y ACCESIBILIDAD
% ==============================================================================
\section{Gestión de Calidad (Durante el Desarrollo y al Final, Herramientas)}

Se revisará completamente todo el desarrollo al final de cada iteración y se corregirán los aspectos necesarios para alcanzar nuestros objetivos, asegurando una detección temprana de errores y validación continua del trabajo entregado.

\subsection{Procedimiento de Calidad en Cada Iteración}
Por cada iteración seguiremos rigurosamente los siguientes pasos:
\begin{enumerate}[leftmargin=*]
  \item \textbf{Definición de criterios de aceptación:} Discutiremos sobre lo que se ha de realizar en la iteración y qué criterios seguimos para determinar si se ha completado correctamente (Definition of Ready y Definition of Done).
  \item \textbf{Revisiones de código:} Se revisará el código obligatoriamente en busca de errores antes de añadirlo a la rama de desarrollo de la nueva iteración. Se usarán herramientas de debug e  IA para logar la revisión temprana.
  \item \textbf{Análisis Estático Continuo:} Escaneo automatizado del código fuente en cada ciclo para identificar vulnerabilidades, duplicidades y fallos de estilo.
  \item \textbf{Pruebas Unitarias y de Integración:} Ejecución de pruebas automatizadas sobre los nuevos módulos desarrollados en la iteración correspondiente. Este paso puede llegar a ser omitido en las primeras iteraciones debido a la falta de contenido para la realización de pruebas.
  \item \textbf{Demostración y Validación de la Iteración:} Al finalizar cada iteración, se realizará una revisión funcional con el cliente para presentar y validar el trabajo desarrollado. Esta revisión permitirá comprobar que los resultados obtenidos satisfacen sus necesidades y expectativas.
\end{enumerate}

\subsection{Proceso de Hardening y Verificación Global}
Al terminar la tercera iteración, se llevará a cabo el \textbf{proceso de hardening y verificación global}, el cual constará de:
\begin{itemize}[leftmargin=*]
  \item Una revisión funcional integral por parte del cliente.
  \item Una auditoría final exhaustiva que compruebe el correcto funcionamiento de la aplicacion.
\end{itemize}

\subsection{Cuadro Oficial de Herramientas de Calidad}

\vspace{0.2cm}
\noindent
{\small
\begin{tabularx}{\textwidth}{|>{\raggedright\arraybackslash\bfseries}p{3.8cm}|>{\raggedright\arraybackslash\bfseries}p{3.2cm}|X|}
\hline
\rowcolor{tableHeader}
Fase / Ámbito & Herramienta & Aplicación en el Proyecto \\
\hline
Control de Versiones \& CI/CD & GitHub & Gestión del repositorio, control de Pull Requests y ramas. \\
\hline
Análisis Estático \& Estilo & SonarCloud / ESLint & Detección automática de bugs, vulnerabilidades y formato de código. \\
\hline
Pruebas Unitarias \& API & Jest & Automatización de pruebas de lógica e integración de servicios. \\
\hline
Accesibilidad \& UX & Lighthouse / WAVE & Evaluación de niveles de accesibilidad (contraste, etiquetas, lectura). \\
\hline
Gestión de Tareas y Bugs & Jira & Trazabilidad de requisitos, incidencias y feedback de las demos. \\
\hline
\end{tabularx}
}
\vspace{0.3cm}
